% !TEX program = xelatex

\documentclass[12pt]{exam}
\AtBeginDocument{\Large}
\usepackage[margin=1in]{geometry}
\usepackage{xeCJK}
\usepackage{amsmath,amssymb}
\usepackage{comment}
\excludecomment{answer}
\setCJKmainfont{Noto Serif CJK SC}

\title{\vspace{-2cm}椭圆练习题\vspace{-1cm}}
\author{}
\date{}

\begin{document}
\maketitle
\vspace{-1.5cm}

\begin{questions}
\question 已知椭圆的方程为 $x^2 + 2y^2 = 2$。过该椭圆的焦点引斜率为 $1$ 的直线 $l$，直线 $l$ 与椭圆相交于 $A, B$ 两点。求由点 $A, B$ 以及椭圆中心 $O$ 所构成的 $\triangle OAB$ 的面积。

\end{questions}

\begin{answer}
\section*{解答}
将椭圆方程化为标准方程：
\[ \frac{x^2}{2} + y^2 = 1 \]
由此可知，$a^2 = 2$，$b^2 = 1$。
椭圆的半焦距 $c = \sqrt{a^2 - b^2} = \sqrt{2 - 1} = 1$。
因此，椭圆的焦点坐标为 $F_1(-1, 0)$ 和 $F_2(1, 0)$。

由于图形的对称性，无论选择哪一个焦点，所构成的三角形面积均相同。不妨取右焦点 $F(1, 0)$。

已知直线 $l$ 过点 $F(1, 0)$ 且斜率 $k = 1$，由点斜式可得直线 $l$ 的方程为：
\[ y - 0 = 1 \cdot (x - 1) \implies y = x - 1 \]

将直线方程 $y = x - 1$ 代入椭圆方程 $x^2 + 2y^2 = 2$ 中，得：
\[ x^2 + 2(x - 1)^2 = 2 \]
展开并整理得：
\[ x^2 + 2(x^2 - 2x + 1) = 2 \]
\[ 3x^2 - 4x = 0 \]
解得：
\[ x_1 = 0, \quad x_2 = \frac{4}{3} \]

将 $x$ 的值代回直线方程 $y = x - 1$，得到对应的纵坐标：
当 $x_1 = 0$ 时，$y_1 = -1$，即交点 $A(0, -1)$；\\
当 $x_2 = \frac{4}{3}$ 时，$y_2 = \frac{1}{3}$，即交点 $B\left(\frac{4}{3}, \frac{1}{3}\right)$。

椭圆的中心为原点 $O(0, 0)$。
由两点间距离公式，弦长 $|AB|$ 为：
\[ |AB| = \sqrt{\left(\frac{4}{3} - 0\right)^2 + \left(\frac{1}{3} - (-1)\right)^2} = \sqrt{\frac{16}{9} + \frac{16}{9}} = \frac{4\sqrt{2}}{3} \]

原点 $O(0, 0)$ 到直线 $l: x - y - 1 = 0$ 的距离（即 $\triangle OAB$ 在 $AB$ 边上的高）$d$ 为：
\[ d = \frac{|0 - 0 - 1|}{\sqrt{1^2 + (-1)^2}} = \frac{1}{\sqrt{2}} = \frac{\sqrt{2}}{2} \]

因此，$\triangle OAB$ 的面积为：
\[ S_{\triangle OAB} = \frac{1}{2} |AB| \cdot d = \frac{1}{2} \times \frac{4\sqrt{2}}{3} \times \frac{\sqrt{2}}{2} = \frac{2}{3} \]

\textit{注：若选取左焦点 $F_1(-1, 0)$，直线方程为 $y = x + 1$，同理可求得交点为 $(0, 1)$ 和 $\left(-\frac{4}{3}, -\frac{1}{3}\right)$，计算出的面积同样为 $\frac{2}{3}$。}

\vspace{0.5cm}
\textbf{答案：} $\triangle OAB$ 的面积为 $\frac{2}{3}$。

\end{answer}

\begin{questions}
\question 已知椭圆方程为 $\frac{(x-2)^2}{4} + (y-1)^2 = 1$。定抛物线方程为 $y = \frac{5}{16}(x-2)^2 - \frac{1}{4}$。问椭圆与抛物线的交点个数为多少？
\end{questions}

\begin{answer}
首先，将抛物线方程 $y = \frac{5}{16}(x-2)^2 - \frac{1}{4}$ 进行变形，用 $y$ 来表示 $(x-2)^2$：
\[ (x-2)^2 = \frac{16}{5}\left(y + \frac{1}{4}\right) \]
\[ (x-2)^2 = \frac{16}{5}y + \frac{4}{5} \]

接着，将椭圆方程 $\frac{(x-2)^2}{4} + (y-1)^2 = 1$ 等式两边同乘 $4$，展开得到：
\[ (x-2)^2 + 4(y-1)^2 = 4 \]

将抛物线变形后的表达式代入上述椭圆方程中，消去 $x$：
\[ \left(\frac{16}{5}y + \frac{4}{5}\right) + 4(y^2 - 2y + 1) = 4 \]

展开并整理该等式：
\[ \frac{16}{5}y + \frac{4}{5} + 4y^2 - 8y + 4 = 4 \]
\[ 4y^2 - \frac{24}{5}y + \frac{4}{5} = 0 \]

为了方便计算，将等式两边同除以 $4$ 并乘以 $5$，转化为整系数方程：
\[ 5y^2 - 6y + 1 = 0 \]

对该一元二次方程进行因式分解：
\[ (5y - 1)(y - 1) = 0 \]

解得 $y$ 的两个可能值为：
\[ y_1 = 1 \quad \text{或} \quad y_2 = \frac{1}{5} \]

最后，将求得的 $y$ 值代回抛物线变形后的方程 $(x-2)^2 = \frac{16}{5}y + \frac{4}{5}$ 中，求解对应的 $x$ 值：

\begin{itemize}
    \item \textbf{当 $y = 1$ 时：}
    \[ (x-2)^2 = \frac{16}{5}(1) + \frac{4}{5} = \frac{20}{5} = 4 \]
    解得 $x - 2 = \pm 2$，即 $x = 4$ 或 $x = 0$。\\
    由此得到两个交点坐标：$(4, 1)$ 和 $(0, 1)$。

    \item \textbf{当 $y = \frac{1}{5}$ 时：}
    \[ (x-2)^2 = \frac{16}{5}\left(\frac{1}{5}\right) + \frac{4}{5} = \frac{16}{25} + \frac{20}{25} = \frac{36}{25} \]
    解得 $x - 2 = \pm \frac{6}{5}$，即 $x = \frac{16}{5}$ 或 $x = \frac{4}{5}$。\\
    由此又得到两个交点坐标：$\left(\frac{16}{5}, \frac{1}{5}\right)$ 和 $\left(\frac{4}{5}, \frac{1}{5}\right)$。
\end{itemize}

综上所述，方程组产生了 $4$ 组完全不同的实数解。

\textbf{最终答案：} 交点个数为 $4$ 个。
\end{answer}


\end{document}
